% Source: normalization_in_cnn.tex
\begin{notechapter}
% ---- chapter-local definitions ----
\lstdefinestyle{python}{
    language=Python,
    basicstyle=\ttfamily\small,
    keywordstyle=\color{blue},
    commentstyle=\color{green!50!black},
    stringstyle=\color{red},
    showstringspaces=false,
    breaklines=true,
    frame=single,
    numbers=left,
    numberstyle=\tiny\color{gray},
    captionpos=b
}
% -----------------------------------
\chapter{Normalization in Convolutional Networks}\label{ch:cnn-normalization}

\section{Introduction}
Normalization is a crucial preprocessing step when working with
Convolutional Neural Networks (CNNs). It involves scaling and
standardizing input data (e.g., images) to improve the training
process and model performance. This document explains the
importance of normalization, common techniques, and how to compute
and apply per-channel normalization.

\section{Why Normalize Images?}
Normalizing images provides several benefits:
\begin{itemize}
    \item \textbf{Faster Convergence}: Normalized inputs help
          optimization algorithms (e.g., gradient descent) converge
          faster by smoothing the loss landscape.
    \item \textbf{Numerical Stability}: Normalization prevents issues
          like exploding or vanishing gradients caused by large or
          small pixel values.
    \item \textbf{Consistent Feature Scaling}: Ensures all input
          features (pixel values) are on a similar scale, which is
          important for gradient-based optimization.
    \item \textbf{Better Generalization}: Normalized inputs can help
          the model generalize better to unseen data.
\end{itemize}

\section{Common Normalization Techniques}
There are several ways to normalize images for CNNs:

\subsection{Scaling to [0, 1]}
Divide each pixel value by the maximum possible value (e.g., 255 for
8-bit images):
\[
x_{\text{normalized}} = \frac{x}{255}
\]
This is the simplest and most common normalization technique.

\subsection{Scaling to [-1, 1]}
Rescale pixel values to the range [-1, 1] by subtracting 127.5 and
then dividing by 127.5:
\[
x_{\text{normalized}} = \frac{x - 127.5}{127.5}
\]
This is often used in models like GANs (Generative Adversarial
Networks).

\subsection{Standardization (Zero Mean, Unit Variance)}
Subtract the mean pixel value and divide by the standard deviation:
\[
x_{\text{standardized}} = \frac{x - \mu}{\sigma}
\]
Here, \(\mu\) is the mean pixel value, and \(\sigma\) is the
standard deviation of pixel values. This is particularly useful
when the dataset has a known mean and standard deviation (e.g.,
ImageNet).

\subsection{Per-Channel Normalization}
Per-channel normalization is a specific type of standardization
where the mean and standard deviation are computed \textbf{separately
for each channel} (e.g., Red, Green, Blue for RGB images). This is
commonly used in pre-trained models like ResNet, VGG, and Inception,
which were trained on ImageNet\footnote{ImageNet is a large-scale
dataset of over 14 million labeled images, categorized into more
than 20,000 classes. It is widely used for training and benchmarking
deep learning models, particularly in computer vision. The ImageNet
Large Scale Visual Recognition Challenge (ILSVRC) has driven
significant advancements in image classification, object detection,
and other vision tasks. Many pre-trained models (e.g., VGG, ResNet,
Inception) are trained on ImageNet and achieve state-of-the-art
performance.} with per-channel normalization.

\section{Standardization vs. Per-Channel Normalization}
Standardization and per-channel normalization are closely related but
not exactly the same. Here’s the distinction:

\subsection{Standardization}
Standardization refers to transforming data to have a \textbf{mean of 0}
and a \textbf{standard deviation of 1}. The formula is:
\[
x_{\text{standardized}} = \frac{x - \mu}{\sigma}
\]
where \(\mu\) is the mean and \(\sigma\) is the standard deviation.
Standardization can be applied globally (to the entire dataset) or
per-channel (to each channel separately).

\subsection{Per-Channel Normalization}
Per-channel normalization is a specific type of standardization where
the mean and standard deviation are computed \textbf{separately for
each channel} (e.g., Red, Green, Blue for RGB images). For example:
\[
x_{\text{normalized}}^{(R)} = \frac{x^{(R)} - \mu_R}{\sigma_R}, \quad
x_{\text{normalized}}^{(G)} = \frac{x^{(G)} - \mu_G}{\sigma_G}, \quad
x_{\text{normalized}}^{(B)} = \frac{x^{(B)} - \mu_B}{\sigma_B}
\]
Per-channel normalization is commonly used in pre-trained models like
ResNet, VGG, and Inception, which were trained on ImageNet with
per-channel normalization.

\subsection{Key Differences}
\begin{itemize}
    \item \textbf{Scope}: Standardization can be applied globally or
          per-channel, while per-channel normalization is specific to
          multi-channel data like images.
    \item \textbf{Application}: Standardization is a general technique
          for any type of data, while per-channel normalization is
          specific to images.
\end{itemize}

\subsection{When Are They the Same?}
If you apply standardization \textbf{per-channel} to an image, then
per-channel normalization and standardization are effectively the
same thing. However, if you apply standardization \textbf{globally},
then it is not the same as per-channel normalization.

\section{Per-Channel Normalization in Pre-Trained Models}
Many pre-trained models, such as ResNet, VGG, and Inception, use
per-channel normalization because they were trained on ImageNet with
this technique. When using these models, it is important to apply
the same normalization parameters to your data to ensure compatibility
with the pre-trained weights.

\subsection{Why Compute Mean and Std on the Training Set Only?}
\begin{itemize}
    \item \textbf{Avoid Data Leakage}: The validation and test sets
          should remain unseen during training to ensure a fair
          evaluation of the model's performance. Computing
          normalization parameters on the entire dataset would leak
          information from the validation and test sets into the
          training process.
    \item \textbf{Simulate Real-World Scenarios}: In real-world
          applications, the model will encounter new, unseen data.
          Normalizing the training data independently ensures that
          the model is trained and evaluated under conditions that
          mimic real-world deployment.
\end{itemize}

\subsection{How to Handle Normalization for Validation and Test Sets}
\begin{itemize}
    \item \textbf{Use Training Set Parameters}: After computing the
          mean and standard deviation from the training set, apply
          the same normalization parameters to the validation and
          test sets.
    \item \textbf{No Recalculation}: Do not recompute the mean and
          standard deviation for the validation or test sets, as
          this would introduce bias.
\end{itemize}

\subsection{Example: Computing Mean and Std in Python}
Below is an example of computing the mean and standard deviation
for the CIFAR-10 training set:

\begin{lstlisting}[language=Python, style=python]
from torchvision import datasets, transforms
from torch.utils.data import DataLoader
import numpy as np

# Load the training set
train_dataset = datasets.CIFAR10(root='./data', train=True,
                                 download=True, transform=transforms.ToTensor())

# Compute mean and std from the training set
def compute_mean_std(dataset):
    loader = DataLoader(dataset, batch_size=32, shuffle=False)
    mean = 0.0
    std = 0.0
    total_images = 0

    for images, _ in loader:
        batch_size, channels, height, width = images.shape
        images = images.view(batch_size, channels, -1)
        mean += images.mean(2).sum(0)
        std += images.std(2).sum(0)
        total_images += batch_size

    mean /= total_images
    std /= total_images
    return mean, std

mean, std = compute_mean_std(train_dataset)
print(f"Mean: {mean}, Std: {std}")

# Define normalization transform using training set parameters
transform = transforms.Compose([
    transforms.ToTensor(),
    transforms.Normalize(mean, std)
])

# Apply the same transform to validation and test sets
val_dataset = datasets.CIFAR10(root='./data', train=False,
                               download=True, transform=transform)
test_dataset = datasets.CIFAR10(root='./data', train=False,
                                download=True, transform=transform)
\end{lstlisting}

\section{Key Takeaways}
\begin{itemize}
    \item Normalizing images is a best practice when training CNNs.
    \item Common normalization techniques include scaling to [0, 1],
          scaling to [-1, 1], and standardization (zero mean, unit
          variance).
    \item Per-channel normalization is a specific type of
          standardization used in pre-trained models like ResNet,
          VGG, and Inception.
    \item Compute normalization parameters on the training set only
          to avoid data leakage.
    \item Apply the same normalization parameters to the validation
          and test sets for consistency.
\end{itemize}
\end{notechapter}
